\documentclass[11pt,a4paper]{article}

% ── Codificación y fuentes ────────────────────────────────────────────────────
\usepackage[utf8]{inputenc}
\usepackage[T1]{fontenc}
\usepackage{lmodern}
\usepackage[spanish]{babel}

% ── Geometría y espaciado ─────────────────────────────────────────────────────
\usepackage[top=2.5cm, bottom=2.5cm, left=2.8cm, right=2.8cm]{geometry}
\usepackage{setspace}
\setstretch{1.15}
\usepackage{parskip}

% ── Tablas ────────────────────────────────────────────────────────────────────
\usepackage{booktabs}
\usepackage{tabularx}
\usepackage{array}
\usepackage{longtable}
\usepackage{enumitem}

% ── Colores y cajas ───────────────────────────────────────────────────────────
\usepackage[dvipsnames]{xcolor}
\usepackage{tcolorbox}
\tcbuselibrary{skins, breakable}

% ── Cabeceras y pies de página ────────────────────────────────────────────────
\usepackage{fancyhdr}
\usepackage{lastpage}
\pagestyle{fancy}
\fancyhf{}
\fancyhead[L]{\small\textbf{Informe de Evaluación} --- Física para Bioanalistas}
\fancyhead[R]{\small debora echezuria 32585229}
\fancyfoot[C]{\small Página \thepage\ de \pageref{LastPage}}
\renewcommand{\headrulewidth}{0.4pt}

% ── Hiperenlaces ──────────────────────────────────────────────────────────────
\usepackage[colorlinks=true, linkcolor=NavyBlue, urlcolor=NavyBlue]{hyperref}

% ── Paleta de colores personalizados ─────────────────────────────────────────
\definecolor{desempeno_excelente}{HTML}{1A6B2F}
\definecolor{desempeno_bueno}{HTML}{2E5A9C}
\definecolor{desempeno_suficiente}{HTML}{8B6914}
\definecolor{desempeno_deficiente}{HTML}{C0392B}
\definecolor{desempeno_insuficiente}{HTML}{7B241C}
\definecolor{grisFondo}{HTML}{F5F5F5}
\definecolor{azulTitulo}{HTML}{1B2A6B}
\definecolor{verdeFortaleza}{HTML}{1A6B2F}
\definecolor{naranjaMejora}{HTML}{A04000}

% ── Estilos de cajas ─────────────────────────────────────────────────────────
\tcbset{
  cajaTitulo/.style={
    enhanced, breakable,
    colback=azulTitulo!8, colframe=azulTitulo,
    fonttitle=\bfseries\large, coltitle=white,
    attach boxed title to top left={yshift=-2mm, xshift=4mm},
    boxed title style={colback=azulTitulo},
    top=6pt, bottom=6pt, left=8pt, right=8pt,
  },
  cajaFortaleza/.style={
    enhanced, breakable,
    colback=verdeFortaleza!6, colframe=verdeFortaleza!60,
    fonttitle=\bfseries, coltitle=verdeFortaleza,
    top=4pt, bottom=4pt, left=8pt, right=8pt,
  },
  cajaMejora/.style={
    enhanced, breakable,
    colback=naranjaMejora!6, colframe=naranjaMejora!60,
    fonttitle=\bfseries, coltitle=naranjaMejora,
    top=4pt, bottom=4pt, left=8pt, right=8pt,
  },
  cajaRecomendacion/.style={
    enhanced, breakable,
    colback=grisFondo, colframe=gray!50,
    fonttitle=\bfseries, coltitle=black,
    top=4pt, bottom=4pt, left=8pt, right=8pt,
  },
}

% ─────────────────────────────────────────────────────────────────────────────
\begin{document}

% ══ PORTADA ══════════════════════════════════════════════════════════════════
\begin{center}
  {\Large\bfseries\color{azulTitulo} Universidad de Oriente/Núcleo Sucre --- Física para Bioanalistas}\\[4pt]
  {\large Informe de Evaluación Preliminar}\\[12pt]
  \rule{\linewidth}{1.2pt}\\[6pt]
  {\LARGE\bfseries debora echezuria 32585229}\\[4pt]
  \rule{\linewidth}{0.6pt}
\end{center}

% ══ FICHA TÉCNICA ════════════════════════════════════════════════════════════
\vspace{4pt}
\begin{tcolorbox}[cajaTitulo, title=Datos del informe]
\begin{tabular}{@{}llll@{}}
  \textbf{Estudiante:}  & debora echezuria 32585229  &
  \textbf{Fecha:}       & 2026-06-25 \\[4pt]
  \textbf{Archivo PDF:} & \texttt{debora\_echezuria\_32585229.pdf}  &
  \textbf{Páginas:}     & 4 \\
\end{tabular}
\end{tcolorbox}

% ══ RESULTADO GLOBAL ═════════════════════════════════════════════════════════
\vspace{6pt}
\begin{center}
  \begin{tcolorbox}[
    enhanced,
    colback=white, colframe=desempeno_suficiente,
    width=0.62\linewidth,
    boxrule=2pt,
    halign=center, valign=center,
    top=8pt, bottom=8pt,
  ]
    \centering
    {\large\bfseries Puntaje sugerido}\\[6pt]
    {\Huge\bfseries\color{desempeno_suficiente} 6.0/10}\\[6pt]
    {\large Nivel de desempeño: \textbf{\color{desempeno_suficiente} Suficiente}}
  \end{tcolorbox}
\end{center}

% ══ RESUMEN GENERAL ══════════════════════════════════════════════════════════
\section*{Resumen general}
El estudiante demuestra un buen dominio de los principios básicos de la hidrostática y el principio de Arquímedes, aplicando correctamente las fórmulas y realizando los cálculos. Sin embargo, presenta deficiencias significativas en la aplicación de conceptos de hemodinámica más complejos, como las ecuaciones de continuidad y Bernoulli, y la ley de Poiseuille, donde se observan omisiones, errores conceptuales en la interpretación de resultados o errores procedimentales en los cálculos.

% ══ FORTALEZAS Y ÁREAS DE MEJORA ════════════════════════════════════════════
\vspace{4pt}
\begin{minipage}[t]{0.48\linewidth}
  \begin{tcolorbox}[cajaFortaleza, title={Fortalezas}]
\begin{itemize}[leftmargin=1.5em, itemsep=2pt]
  \item Correcta aplicación de los principios de presión hidrostática y Arquímedes.
  \item Habilidad para extraer datos relevantes y realizar conversiones de unidades de manera precisa en la mayoría de los casos.
  \item Presentación clara y organizada de los pasos para las preguntas resueltas correctamente.
\end{itemize}
  \end{tcolorbox}
\end{minipage}
\hfill
\begin{minipage}[t]{0.48\linewidth}
  \begin{tcolorbox}[cajaMejora, title={Areas de mejora}]
\begin{itemize}[leftmargin=1.5em, itemsep=2pt]
  \item Aplicación de las ecuaciones de continuidad y Bernoulli para problemas de flujo.
  \item Interpretación precisa de preguntas que solicitan 'cambio porcentual'.
  \item Precisión en cálculos que involucran potencias y notación científica.
\end{itemize}
  \end{tcolorbox}
\end{minipage}

% ══ OBSERVACIONES POR PREGUNTA ═══════════════════════════════════════════════
\section*{Observaciones por pregunta}

\begin{longtable}{>{\bfseries}p{2.8cm} p{1.6cm} p{7.0cm} p{2.8cm}}
  \toprule
  \textbf{Pregunta} & \textbf{Puntaje} & \textbf{Evaluación} & \textbf{Errores detectados} \\
  \midrule
  \endhead
  \bottomrule
  \endfoot
    \textbf{Pregunta 1: Presión Hidrostática y Venosa} & 2.0/2.0 & La respuesta es completamente correcta. El estudiante identificó la fórmula adecuada, extrajo los datos correctamente, realizó el cálculo de manera precisa y utilizó las unidades consistentes. & \textit{---} \\
    \midrule
    \textbf{Pregunta 2: Principio de Arquímedes} & 2.0/2.0 & La respuesta es completamente correcta. El estudiante calculó la fuerza de empuje, el volumen y la densidad del tejido de forma impecable, aplicando las fórmulas y unidades adecuadas en cada paso. & \textit{---} \\
    \midrule
    \textbf{Pregunta 3: Hemodinámica (Continuidad y Bernoulli)} & 0.0/2.0 & El estudiante solo listó los datos del problema, pero no realizó ningún intento de aplicar las ecuaciones de continuidad y Bernoulli para resolver la pregunta y calcular la presión en la zona de constricción. & \textit{No se presenta solución al problema} \\
    \midrule
    \textbf{Pregunta 4: Ley de Poiseuille (Análisis Proporcional)} & 1.0/2.0 & El estudiante aplicó correctamente la relación proporcional de la Ley de Poiseuille (Q $\propto$ r$^{4}$) y calculó el factor de reducción del flujo (0.4096). Sin embargo, interpretó este resultado como el porcentaje de flujo final (40.96\%) en lugar del 'cambio porcentual' solicitado, que sería una reducción del 59.04\%. & \textit{Error conceptual en la interpretación de 'cambio porcentual'} \\
    \midrule
    \textbf{Pregunta 5: Flujo Viscoso y Resistencia} & 1.0/2.0 & El estudiante identificó correctamente la Ley de Poiseuille y la despejó para la diferencia de presión. El numerador de la ecuación fue calculado correctamente. Sin embargo, se cometió un error procedimental significativo al calcular el término r$^{4}$ en el denominador (el exponente de la potencia de 10 es incorrecto), lo que llevó a un resultado final numéricamente erróneo (aproximadamente un factor de 10). & \textit{Error procedimental en el cálculo de (4 x 10$^{-}$$^{4}$)$^{4}$ (error de exponente)} \\
    \midrule
\end{longtable}

% ══ ERRORES DETECTADOS ═══════════════════════════════════════════════════════
\section*{Análisis de errores}

\subsection*{Errores conceptuales}
\begin{itemize}[leftmargin=1.5em, itemsep=2pt]
  \item Confusión entre el porcentaje de un valor final y el porcentaje de cambio (Pregunta 4).
  \item Falta de aplicación de los principios de continuidad y Bernoulli para resolver un problema de hemodinámica (Pregunta 3).
\end{itemize}

\subsection*{Errores procedimentales}
\begin{itemize}[leftmargin=1.5em, itemsep=2pt]
  \item Error en el cálculo de una potencia con notación científica, específicamente (4 x 10$^{-}$$^{4}$)$^{4}$, afectando el exponente final (Pregunta 5).
\end{itemize}

% ══ RECOMENDACIONES AL ESTUDIANTE ════════════════════════════════════════════
\section*{Retroalimentación para el estudiante}
\begin{tcolorbox}[cajaRecomendacion, title={Dirigido a: debora echezuria 32585229}]
Débora, has demostrado una base sólida en los principios fundamentales de la hidrostática y Arquímedes. Para mejorar, te sugiero que dediques más tiempo a practicar problemas de dinámica de fluidos, especialmente aquellos que involucran las ecuaciones de continuidad, Bernoulli y Poiseuille. Es crucial que no dejes preguntas en blanco y que intentes aplicar los conceptos aprendidos. Presta especial atención a la lectura de las preguntas para interpretar correctamente lo que se solicita (por ejemplo, 'cambio porcentual' vs. 'porcentaje final'). Finalmente, revisa tus cálculos, sobre todo aquellos con potencias y notación científica, para evitar errores que afecten el resultado final.
\end{tcolorbox}

% ══ NOTA PARA EL PROFESOR ════════════════════════════════════════════════════
\section*{Nota para el profesor}
\begin{tcolorbox}[
  enhanced, breakable,
  colback=yellow!8, colframe=yellow!60!black,
  fonttitle=\bfseries, coltitle=black,
  title={Observaciones para revisión manual},
  top=4pt, bottom=4pt, left=8pt, right=8pt,
]
La estudiante Débora Echeverría muestra un buen desempeño en los primeros dos problemas. El problema 3 fue dejado en blanco. En el problema 4, la estudiante calculó correctamente la relación de flujos pero no el cambio porcentual. En el problema 5, hubo un error de cálculo en el denominador (exponente de r$^{4}$), lo que resultó en un valor final incorrecto por un factor de 10. El puntaje sugerido refleja estos puntos.
\end{tcolorbox}

% ══ PIE DE INFORME ════════════════════════════════════════════════════════════
\vspace{12pt}
\begin{center}
  \footnotesize\color{gray}
  Informe generado automáticamente por el sistema de evaluación con IA (gemini-2.5-flash) y soporte LaTex. \\
  Este documento es una evaluación \textbf{preliminar} para ser revisado por el profesor antes de ser entregado al 
  estudiante. \\
  Generado el 2026-06-25 \textbullet\ BIOANALISIS Sistema Académico de Evaluación.
  \end{center}

\end{document}
